\section{Introduction}
\IEEEPARstart{C}{onversational} business-intelligence interfaces let non-technical users question operational data in natural language~\cite{setlur2016eviza,fast2018iris,quamar2022nli}. For relational databases, Text-to-SQL has matured around large benchmarks such as Spider and BIRD~\cite{spider2018,bird2023}. Document stores are harder targets. In MongoDB the query is an aggregation pipeline whose stages reshape the document stream, arrays must be unwound before grouping, and the operative schema is only partly visible in the stored documents~\cite{mongodb2024manual,tend2025benchmark}. A recent MongoDB-native benchmark reports that the best system it evaluates remains below 36\% mean execution accuracy~\cite{tend2025benchmark}, and earlier NoSQL translators cover basic operations such as find, insert, update and remove~\cite{hossen2023bert}.

These evaluations pair every request with one intended answer. Real requests are often underspecified. ``Show the top products'' does not say whether \emph{top} refers to price, quantity or another numeric field. A single-pass generator must choose, and because the resulting pipeline executes and returns plausible rows, the user cannot tell that a choice was made. We call this \emph{silent assumption}. Clarification questions are established in conversational search and open-domain question answering~\cite{aliannejadi2019asking,zamani2020generating,min2020ambigqa}; this paper studies a deterministic, database-grounded variant for MongoDB.

AMB-MQL places a lightweight gate before generation. After profiling a five-document sample of each in-scope collection, the gate asks for clarification when the query contains one of 15 vague terms, the scope offers at least two numeric candidate fields, and the query does not already bind one. It returns up to four candidate fields plus an ``other'' option. The user's selection is validated against the offered options and the schema, merged into the prompt, and a single model call generates the pipeline. A stage allowlist and a collection-scope check validate the pipeline before execution, and mutating pipelines require a single-use signed approval token. The service is multi-tenant: connection strings are encrypted at rest and database handles are pooled per connection.

We evaluate the gate and the end-to-end flow on three MongoDB sample datasets with \numTotalQueries{} queries labelled by \numAnnotators{}. Against ungated baselines, the gate reduces silent assumptions on lexically cued ambiguous queries, adds nothing on paraphrased ambiguity where the language model's own clarification mode already asks, and triggers unnecessary clarifications on clear queries; its four-candidate list can also omit the intended metric. The study covers metric ambiguity only, one model, and a simulated user.

The contributions are:
\begin{itemize}
\item A deterministic, schema-grounded ambiguity gate with allowlist-validated clarification cards, and a precise account of what it detects and misses.
\item An end-to-end MongoDB service architecture: profiling, hybrid clarification, single-call generation, stage-allowlist validation, human-in-the-loop write staging, tenant credential vault and connection pool.
\item A human-labelled benchmark of \numTotalQueries{} scoped queries over three Atlas sample datasets, with metrics that separate lexical, paraphrase and structural ambiguity.
\end{itemize}
